\documentclass[twoside]{article}

% AISTATS 2027 author kit was not yet available when this draft was written.
% Use the latest official style without modification and replace it when the
% 2027 kit is released.
\usepackage{aistats2026}
\usepackage[round]{natbib}
\renewcommand{\bibname}{References}
\renewcommand{\bibsection}{\subsubsection*{\bibname}}
\bibliographystyle{apalike}

\usepackage{amsmath,amssymb,amsthm,mathtools,bm}
\usepackage{booktabs,array,multirow}
\usepackage[hidelinks]{hyperref}

\newcommand{\E}{\mathbb{E}}
\newcommand{\Pp}{\mathbb{P}}
\newcommand{\Var}{\operatorname{Var}}
\newcommand{\Cov}{\operatorname{Cov}}
\newcommand{\tr}{\operatorname{tr}}
\newcommand{\cF}{\mathcal{F}}
\newcommand{\cR}{\mathcal{R}}
\newcommand{\cS}{\mathcal{S}}
\newcommand{\cX}{\mathcal{X}}
\newcommand{\pos}[1]{[#1]_{+}}
\newcommand{\TBD}{\textsc{pending}}

\theoremstyle{plain}
\newtheorem{theorem}{Theorem}
\newtheorem{proposition}{Proposition}
\newtheorem{corollary}{Corollary}
\newtheorem{lemma}{Lemma}
\theoremstyle{definition}
\newtheorem{assumption}{Assumption}
\theoremstyle{remark}
\newtheorem{remark}{Remark}

\begin{document}

\runningtitle{When Is One Flow Enough?}
\runningauthor{Anonymous Authors}

\twocolumn[
\aistatstitle{When Is One Flow Enough?\\
Learning Damage and Restoration as Opposing Neural Flows}
\aistatsauthor{Anonymous Author(s)}
\aistatsaddress{Anonymous Institution}
]

\begin{abstract}
A bounded trajectory records only a net change even when two nonnegative
processes act simultaneously. We study when a single signed flow is sufficient
and when damage/interruption and restoration should be learned separately. For
\(m_0(x,y)=D_0(x)(1-y)-R_0(x)y\), we derive the exact population loss of the
best state-scaled one-flow approximation,
\[
G(x)=v(x)\min\!\left\{\frac{R_0(x)^2}{A(x)},
                       \frac{D_0(x)^2}{B(x)}\right\},
\]
where \(A=\E[(1-Y)^2\mid X=x]\), \(B=\E[Y^2\mid X=x]\), and
\(v=\Var(Y\mid X=x)\). Thus one flow is sufficient exactly when one
physical flow is absent or the conditional state is degenerate. An orthogonal
reparameterization separates net drift from total turnover and yields the exact
conditional information matrix \(\operatorname{diag}(1,v(x))\). A
fixed-design Gaussian benchmark then isolates the representation--estimation
crossover, while an exactly solvable two-state case varies only state
dispersion and reproduces the predicted threshold. Because these results are
one-step statements, we derive an exact rollout-error recursion and design one
focused real-storm experiment that reports both held-out teacher-forced
transition error and 24-hour recursive forecast error. The proposed two-flow
neural dynamics and a parameter-matched signed one-flow network use identical
inputs, state scaling, optimization, and event-balanced rollout loss; gradient
boosting and damped persistence provide practical reference points. We make no
claim that two flows dominate every event or every horizon.
\end{abstract}

\section{INTRODUCTION}

A scalar state may rise and fall because two processes act at the same time.
During severe weather, some customers lose service while others are restored.
Forecasting models often replace the two nonnegative flows by one signed net
update. This reduces the number of functions to estimate, but it removes the
possibility that both directions are active under the same covariates.

Power-outage evolution makes the model-selection problem concrete. Utility
resilience curves can be decomposed retrospectively into overlapping outage and
restore processes \citep{carrington2021}, while public county records expose
only the aggregate fraction of customers without power. We ask:
\begin{quote}
\emph{When is one signed flow enough, and when is the representation gain from
separating damage and restoration large enough to repay the estimation cost of
a second function?}
\end{quote}

The paper separates three questions. \emph{Necessity}: does collapsing two
flows create irreducible approximation error? \emph{Learnability}: does the
conditional state distribution contain enough information to distinguish the
two rates? \emph{Finite-sample usefulness}: does the representation gain exceed
the additional estimation and misspecification cost? The real-data experiment
then asks whether this one-step structural benefit survives recursive
prediction on a held-out storm.

Our contributions are tightly coupled. First, we give the exact constrained
oracle gap between a two-flow conditional mean and the parameter-matched union
of one-flow rays. Second, we show that the natural net-drift/turnover
coordinates are conditionally orthogonal, making the weak direction and its
information explicit. Third, we connect one-step error to trajectory error and
run one leave-one-event-out neural study. The one-flow comparison is the sole
structural ablation; external baselines establish forecasting context rather
than define the theory.

\section{RELATED WORK}

At fixed \(x\), our transition equation is a two-column varying-coefficient
regression \citep{hastie1993}. We do not claim varying-coefficient regression as
novel; the contribution is the exact geometry of the opposing-flow basis
\((1-y,-y)\) and the resulting model-selection criterion. Neural differential
models provide flexible parameterizations of unknown dynamics
\citep{chen2018}. Our focus is when two directional components are
observationally distinguishable and predictively worth estimating.

Outage and restoration processes overlap in utility records
\citep{carrington2021}. We use public EAGLE-I county outage data
\citep{brelsford2024,tansakul2023} and ERA5 weather drivers
\citep{hersbach2020}. Unlike retrospective process decomposition, our learner
observes only the net state and must decide whether two predictive functions are
justified.

\section{MODEL CLASSES AND SCOPE}

Let \(Y_t\in[0,1]\), \(X_t\in\cX\), and \(Z_t=Y_{t+1}-Y_t\).

\begin{assumption}[Pooled conditional two-flow mean]
\label{ass:mean}
There exist measurable bounded functions \(D_0,R_0:\cX\to[0,\infty)\) such
that
\begin{equation}
\E[Z_t\mid X_t=x,Y_t=y]
=D_0(x)(1-y)-R_0(x)y,
\label{eq:drift}
\end{equation}
and the residual has conditional mean zero and finite variance.
\end{assumption}

Assumption~\ref{ass:mean} is a predictive conditional-mean restriction, not a
causal claim. If unobserved county heterogeneity remains after conditioning on
\(X\), and county composition changes with \(Y\), the pooled mean need not be
affine in \(y\). In that case the learned functions are pooled predictive
components rather than physical or crew-productivity parameters.

At fixed \(x\), the two-flow class is
\begin{equation}
\cF_2(x)=\{d(1-y)-ry:d,r\ge0\}.
\end{equation}
The parameter-matched one-flow class is
\begin{equation}
\cF_1(x)=\{a(1-y):a\ge0\}\cup\{-by:b\ge0\}.
\label{eq:oneclass}
\end{equation}
Equivalently, one signed coefficient may switch direction across \(x\), but at
a fixed \(x\) it cannot keep damage and restoration simultaneously positive.
Define
\begin{align}
A(x)&=\E[(1-Y)^2\mid X=x], & B(x)&=\E[Y^2\mid X=x],\nonumber\\
C(x)&=\E[Y(1-Y)\mid X=x], & v(x)&=\Var(Y\mid X=x).
\label{eq:moments}
\end{align}
A direct expansion gives \(A(x)B(x)-C(x)^2=v(x)\).

\section{WHEN IS ONE FLOW ENOUGH?}

\begin{theorem}[Exact oracle one-flow gap]
\label{thm:gap}
Under Assumption~\ref{ass:mean}, let
\[
G(x)=\inf_{g\in\cF_1(x)}
\E[(D_0(x)(1-Y)-R_0(x)Y-g(Y))^2\mid X=x].
\]
If \(A(x),B(x)>0\), then
\begin{equation}
\boxed{
G(x)=v(x)\min\!\left\{\frac{R_0(x)^2}{A(x)},
                        \frac{D_0(x)^2}{B(x)}\right\}.}
\label{eq:gap}
\end{equation}
If \(A(x)=0\) or \(B(x)=0\), then \(v(x)=0\) and \(G(x)=0\). Therefore
\begin{equation}
G(x)>0\quad\Longleftrightarrow\quad
v(x)>0,\ D_0(x)>0,\ R_0(x)>0.
\label{eq:gapstrict}
\end{equation}
\end{theorem}

The theorem applies to the specific state-scaled one-flow comparator in
Equation~\eqref{eq:oneclass}; it is not a statement about every possible scalar
predictor \(a(x)h(y)\). It separates three zero-gap regimes: damage only,
restoration only, and degenerate state support. Only when both directions are
active and the same driver condition occurs at different states is one flow
representationally insufficient.

\subsection{Orthogonal net-drift and turnover coordinates}

Let \(\mu(x)=\E[Y\mid X=x]\) and define
\begin{equation}
\alpha(x)=(1-\mu(x))D_0(x)-\mu(x)R_0(x),
\qquad \tau(x)=D_0(x)+R_0(x).
\label{eq:alphatau}
\end{equation}
Then
\begin{equation}
m_0(x,y)=\alpha(x)-\tau(x)(y-\mu(x)),
\label{eq:orthogonal}
\end{equation}
and
\(D_0=\alpha+\mu\tau\),
\(R_0=(1-\mu)\tau-\alpha\).

\begin{theorem}[Orthogonal information geometry]
\label{thm:info}
For \(\psi(Y)=(1,-(Y-\mu(x)))^\top\),
\begin{equation}
\E[\psi(Y)\psi(Y)^\top\mid X=x]
=\begin{bmatrix}1&0\\0&v(x)\end{bmatrix}.
\label{eq:diaginfo}
\end{equation}
Equivalently, for \(\phi(Y)=(1-Y,-Y)^\top\) and
\(Q(x)=\E[\phi(Y)\phi(Y)^\top\mid X=x]\),
\begin{equation}
\det Q(x)=v(x),\qquad
v(x)\le\lambda_{\min}(Q(x))\le2v(x),
\label{eq:eig}
\end{equation}
with \(1/2\le\lambda_{\max}(Q(x))\le1\). Hence the two conditional
functions are separately identified from the conditional mean if and only if
\(v(x)>0\).
\end{theorem}

In a homoskedastic fixed design centered at its sample mean,
\begin{equation}
\Var(\widehat\alpha\mid Y)=\frac{\sigma^2}{n},\qquad
\Var(\widehat\tau\mid Y)=\frac{\sigma^2}{n\widehat v},\qquad
\Cov(\widehat\alpha,\widehat\tau\mid Y)=0.
\label{eq:alphatauvar}
\end{equation}
Thus accurate net-drift prediction can coexist with unstable recovery of the
two flows. In the interior fixed-design case, the selection index factors as
\begin{equation}
\Gamma_n=
\underbrace{\frac{n\widehat v}{\sigma^2}}_{I_{\tau,n}}
\underbrace{\min\!\left\{\frac{R_0^2}{\widehat A},
                              \frac{D_0^2}{\widehat B}\right\}}
_{\kappa_{\mathrm{conc},n}}.
\label{eq:factor}
\end{equation}
The first factor is information about total turnover; the second measures how
costly it is to suppress one active direction.

\subsection{A finite-sample benchmark}

Fix states \(y_1,\ldots,y_n\) and observe
\(Z=m+\varepsilon\), \(\varepsilon\sim\mathcal N(0,\sigma^2I_n)\).
Let \(P_2\) project onto the rank-two span generated by \((1-y_i)\) and
\((-y_i)\). Let an oracle choose the better of the two rank-one linear spans,
with projector \(P_1^*\), using the noiseless mean \(m\), and define
\(G_n=n^{-1}\|(I-P_1^*)m\|_2^2\).

\begin{proposition}[Oracle-span bias--variance crossover]
\label{prop:crossover}
If the two-column design has rank two, then
\begin{align}
\E[n^{-1}\|P_2Z-m\|_2^2]&=2\sigma^2/n,\nonumber\\
\E[n^{-1}\|P_1^*Z-m\|_2^2]&=G_n+\sigma^2/n,
\label{eq:risks}
\end{align}
and
\begin{equation}
n^{-1}(\|P_1^*Z-m\|_2^2-\|P_2Z-m\|_2^2)
\overset{d}{=}G_n-\frac{\sigma^2}{n}\chi_1^2.
\label{eq:chisq}
\end{equation}
Thus the two-flow span has lower expected fitted-mean error exactly when
\(nG_n/\sigma^2>1\).
\end{proposition}

This proposition is an oracle-selected, span-relaxed, fixed-design Gaussian
benchmark. The dynamic-panel regressor is predetermined, EAGLE-I errors are
heteroskedastic and include state-measurement error, and neural effective
complexity is not one degree of freedom. The threshold is therefore a calibrated
reference, not a plug-in decision rule for the real data.

\subsection{Exactly solvable two-state case}

Let
\begin{equation}
Y\in\{1/2-\delta,1/2+\delta\},\qquad
\Pp(Y=1/2-\delta)=\Pp(Y=1/2+\delta)=1/2,
\end{equation}
with \(D\ge R>0\). Then
\begin{equation}
G(\delta)=\frac{4R^2\delta^2}{1+4\delta^2},
\qquad \lambda_{\min}(Q)=2\delta^2,
\end{equation}
and an even balanced sample obeys
\begin{equation}
\boxed{
\E[\mathrm{Err}_1-\mathrm{Err}_2]
=\frac{4R^2\delta^2}{1+4\delta^2}-\frac{\sigma^2}{n}.}
\label{eq:solvable}
\end{equation}
The expected-risk crossover is
\begin{equation}
\delta_c=\frac{\sigma}{2\sqrt{nR^2-\sigma^2}},
\end{equation}
when the right-hand side lies in \((0,1/2)\).

The supplied solvable case fixes \(D=0.20\), \(R=0.12\), \(n=40\), and
\(\sigma=0.15\), and varies only \(\delta\). It gives
\(\delta_c=0.1008\). Fifty thousand replications per point match the closed
forms: the maximum standardized deviation of either mean risk is below 1.78,
and the maximum error in the exact win probability is 0.0049. Three reporting
points are shown below.

\begin{table}[t]
\centering
\caption{Exactly solvable two-state case. Risks are fitted-mean MSE.}
\label{tab:solvable}
\small
\begin{tabular}{rrrrr}
\toprule
\(\delta\) & \(\Gamma\) & One-flow & Two-flow & Favored\\
\midrule
0.0494 & 0.247 & 0.000702 & 0.001125 & One\\
0.1008 & 1.000 & 0.001125 & 0.001125 & Tie\\
0.1963 & 3.417 & 0.002485 & 0.001125 & Two\\
\bottomrule
\end{tabular}
\end{table}

This is the paper's only synthetic experiment. It verifies the benchmark and
does not claim an exact neural-network threshold.

\subsection{From one-step error to rollout error}

\begin{proposition}[Rollout-error recursion]
\label{prop:rollout}
Along a fixed driver path, let the true and fitted conditional-mean maps be
\(F_t(y)=D_t+(1-\tau_t)y\) and
\(\widehat F_t(y)=\widehat D_t+(1-\widehat\tau_t)y\). Starting from the same
state, define \(e_t=\widehat y_t-\bar y_t\). Then
\begin{equation}
e_{t+1}=(1-\widehat\tau_t)e_t+\delta_t,
\end{equation}
where
\begin{equation}
\delta_t=(1-\bar y_t)(\widehat D_t-D_t)
-\bar y_t(\widehat R_t-R_t),
\end{equation}
and therefore
\begin{equation}
\boxed{
e_H=\sum_{k=0}^{H-1}\delta_k
\prod_{j=k+1}^{H-1}(1-\widehat\tau_j).}
\label{eq:rollout}
\end{equation}
If \(0\le\widehat\tau_t\le1\), then
\(|e_H|\le\sum_{k<H}|\delta_k|\). A geometric saturation bound additionally
requires a uniform lower bound \(\widehat\tau_t\ge\tau_{\min}>0\).
\end{proposition}

Equation~\eqref{eq:rollout} is the bridge to the real experiment. A lower
average one-step error need not improve every trajectory because signed errors
can cancel or reinforce at different times. Hence the study must report both
held-out teacher-forced transition error and recursive forecast error;
conditional state variance alone is not expected to rank event-level forecast
gains.

\section{NEURAL MODELS}

The proposed model uses two bounded neural rates,
\begin{align}
D_\theta(x)&=c_D\sigma(f_D(x)), &
R_\theta(x)&=c_R\sigma(f_R(x)),\nonumber\\
\widehat Y_{t+1}&=\widehat Y_t+D_\theta(X_t)(1-\widehat Y_t)
-R_\theta(X_t)\widehat Y_t.
\label{eq:twoNN}
\end{align}
Both networks have two ReLU hidden layers of width 32 and caps
\(c_D=c_R=0.25\). The comparator is one width-48 signed network,
\begin{equation}
\widehat Y_{t+1}=\widehat Y_t+
\pos{s_\eta(X_t)}(1-\widehat Y_t)-
\pos{-s_\eta(X_t)}\widehat Y_t.
\label{eq:oneNN}
\end{equation}
The parameter counts differ by less than one percent. Inputs, state scaling,
normalization, optimization, masks, and training budget are identical. Each arm
is initialized by the same flow-matching principle expressed in that arm's own
parameterization; the signed arm is not initialized by subtracting rates
calibrated for the susceptible-flow parameterization.

\section{ONE MAIN REAL-STORM EXPERIMENT}

The empirical question is:
\begin{quote}
\emph{Does allowing simultaneous nonnegative damage and restoration flows
reduce held-out transition error, and does that advantage survive recursive
prediction to 24 hours on a held-out storm event?}
\end{quote}

\subsection{Data, split, and objective}

The cohort is exactly the eleven events in
\texttt{panel\_manifest\_g2-convective-11.json}. No event is selected or
removed according to model performance. Because this cohort informed earlier
exploratory analyses, the study is a prospectively locked re-analysis rather
than an untouched external confirmation; all event-level outcomes are shown.
The physically central event, 2021-06-21, may be used only for an implementation
smoke test and never replaces the eleven-event evaluation.

Each outer fold holds out one complete storm event. Within each of the ten
source events, counties are deterministically split 80/20 for training and
checkpoint validation, with all origins from one county kept on the same side.
Normalization and initialization use source-training rows only. The event is
the inferential unit; seeds quantify optimization variability.

Both neural classes use the same event-balanced 24-hour rollout objective,
\begin{equation}
L_{\mathrm{train}}=\frac1{10}\sum_{e=1}^{10}
L_{\mathrm{roll},e}^{1:24}.
\label{eq:objective}
\end{equation}
A source event is sampled uniformly before drawing a minibatch within it.
Validation is the equal-event mean of the same rollout loss. Teacher-forced
one-step loss is an evaluation endpoint, not a training term. No loss mixture
or coefficient sweep is allowed in the main experiment.

\subsection{Methods and budget}

The main table contains exactly four methods:
\begin{enumerate}
\item the two-flow NN in Equation~\eqref{eq:twoNN};
\item the parameter-matched one-flow NN in Equation~\eqref{eq:oneNN};
\item histogram gradient boosting using the same horizon-available information;
\item damped persistence.
\end{enumerate}
The first comparison is both the theorem-matched contrast and the sole
structural ablation. HGB and persistence measure practical accuracy; they do
not test Theorem~\ref{thm:gap}.

The fit budget is fixed before execution:
\[
2\times11\times3=66\text{ neural fits},\qquad
11\times2=22\text{ HGB fits},
\]
plus negligible persistence fits. No architecture search, memory module,
family analysis, event-shift experiment, information gate, semiparametric term,
or 48-hour campaign starts before the main report is frozen.

\subsection{Hypotheses and inference}

After averaging the three neural seeds within each event, define
\begin{align}
d_e^{\mathrm{step}}
&=\operatorname{MSE}^{\mathrm{TF}}_{\mathrm{one},e}
 -\operatorname{MSE}^{\mathrm{TF}}_{\mathrm{two},e},\nonumber\\
d_e^{24}
&=\operatorname{MSE}^{24}_{\mathrm{one},e}
 -\operatorname{MSE}^{24}_{\mathrm{two},e}.
\label{eq:contrasts}
\end{align}
Positive values favor two flows.

\paragraph{H1: transition-level structural signal.}
The equal-event mean of \(d_e^{\mathrm{step}}\) is positive. This is the
closest held-out predictive counterpart to the oracle gap, but it is not an
estimator of \(G(x)\): the observed current state is noisy, the design is
dynamic, and the neural classes need not attain their population oracles.

\paragraph{H2: 24-hour forecast usefulness.}
The equal-event mean of \(d_e^{24}\) is positive. This asks whether any
transition-level benefit survives finite-sample estimation and the weighted
error propagation in Proposition~\ref{prop:rollout}. H+6 and full-path 24-hour
MSE are secondary diagnostics.

For each endpoint we report all eleven event differences, the equal-event mean
and median, the exact two-sided sign-flip test over all \(2^{11}\) assignments,
a 50,000-resample event bootstrap interval, the positive-event count, and
leave-one-event influence. The inferential gate for each hypothesis is a
positive mean with exact sign-flip \(p<0.05\); the bootstrap interval and sign
count are reported diagnostics, not additional redundant vetoes. The joint
flow-separation claim is an intersection--union claim: it is made only when H1
and H2 both pass at level 0.05.

If H1 passes but H2 does not, the structural signal is lost in rollout. If H2
passes without H1, the forecast gain is reported but not attributed to the
oracle-gap theory. If both fail, the current task lies in a one-flow-sufficient
or estimation/misspecification-dominated regime. HGB may outperform both; that
limits practical competitiveness without changing the mathematical theorem.

\begin{table}[t]
\centering
\caption{Locked main real-storm table. TF-1 is reported only for the two
neural flow classes.}
\label{tab:main}
\small
\begin{tabular}{lrrrr}
\toprule
Method & TF-1 MSE & Path-24 & H+6 & H+24\\
\midrule
Two-flow NN & \TBD & \TBD & \TBD & \TBD\\
One-flow NN & \TBD & \TBD & \TBD & \TBD\\
HGB, same information & -- & -- & \TBD & \TBD\\
Damped persistence & -- & -- & \TBD & \TBD\\
\bottomrule
\end{tabular}
\end{table}

\section{DISCUSSION}

The theory does not claim that two flows always win. The exact gap answers when
the second direction is representationally necessary; the orthogonal geometry
answers when total turnover can be estimated; the fixed-design benchmark shows
how additional variance can overturn representation gain; and the rollout
recursion explains why a one-step benefit can disappear or reverse along a
trajectory.

The independent review identifies four load-bearing limitations. First,
unobserved county heterogeneity can violate pooled affinity. Second, fixed
continuous covariates are approximated through function sharing rather than
literal duplicate observations. Third, the current state is predetermined in a
dynamic panel, so exact fixed-design formulas are benchmarks rather than neural
finite-sample theorems. Fourth, EAGLE-I fractions are heteroskedastic and contain
state-measurement and reporting error; errors in the current state attenuate the
turnover coordinate. These limitations are why the real experiment tests
held-out predictive contrasts rather than a plug-in \(\Gamma\) rule or causal
rate recovery.

The weather path is reanalysis, so the experiment isolates the state-dynamics
model rather than end-to-end weather-forecast uncertainty. The neural model is
also intentionally memoryless beyond the current state. Memory, delayed damage,
and semiparametric structure are plausible later extensions, but introducing
them before resolving the one-flow/two-flow comparison would confound the core
question.

\section{CONCLUSION}

A signed net update is not always a harmless compression of two opposing
processes. One flow is exact when one direction is absent or the conditional
state is degenerate. Otherwise, two flows remove a calculable approximation
gap, but their usefulness depends on information about total turnover and on
the cost of estimating the second function. The solvable case makes this
crossover explicit, and the locked real-storm experiment tests whether it
survives neural estimation and 24-hour rollout.

\bibliography{references}

\clearpage
\appendix
\onecolumn

\section{PROOFS}

\subsection{Proof of Theorem~\ref{thm:gap}}
Fix \(x\), suppress it, and write \(D=D_0(x)\), \(R=R_0(x)\). On the
damage ray,
\[
L_D(a)=A(D-a)^2-2CR(D-a)+BR^2,
\]
with unconstrained minimizer \(a^*=D-RC/A\) and interior minimum
\(R^2(B-C^2/A)=R^2v/A\). The restoration ray analog has
\(b^*=R-DC/B\) and interior minimum \(D^2v/B\).

If \(a^*<0\), then \(D/R<C/A\le\sqrt{B/A}\), so
\(D^2/B<R^2/A\), while
\(DC/B<R C^2/(AB)\le R\), hence \(b^*>0\). Therefore the feasible
restoration branch attains the smaller displayed risk. The symmetric argument
applies if \(b^*<0\). Both unconstrained coefficients cannot be negative when
\(v>0\), because multiplying the resulting strict inequalities would give
\(AB<C^2\). Thus the minimum over the union of nonnegative rays is
Equation~\eqref{eq:gap}. If \(v=0\), then \(Y\) is conditionally constant and
one ray matches the scalar conditional mean exactly. Strict positivity follows
when \(v,D,R\) are all positive.

\subsection{Proof of Theorem~\ref{thm:info}}
Equation~\eqref{eq:diaginfo} follows from
\(\E[Y-\mu\mid X=x]=0\) and \(\E[(Y-\mu)^2\mid X=x]=v\). In the original
basis,
\[
Q=\begin{bmatrix}A&-C\\-C&B\end{bmatrix},
\]
so \(\det Q=AB-C^2=v\). Since
\((1-Y)^2+Y^2\le1\), \(\lambda_{\max}(Q)\le1\). For
\(q=(1,-1)^\top/\sqrt2\), \(q^\top Qq=1/2\), so
\(\lambda_{\max}(Q)\ge1/2\). Therefore
\(\lambda_{\min}=v/\lambda_{\max}\in[v,2v]\). Full rank is equivalent to
\(v>0\).

\subsection{Proof of Proposition~\ref{prop:crossover}}
Because \(\cS_1^*\subset\cS_2\),
\(P_\Delta=P_2-P_1^*\) is a rank-one orthogonal projector. Since
\(m\in\cS_2\), the two-flow error is \(P_2\varepsilon\), with expected
squared norm \(2\sigma^2\). The one-flow error is
\(P_1^*\varepsilon-P_\Delta m\); its orthogonal terms have expected squared
norm \(\sigma^2+nG_n\). Subtraction gives
\[
\mathrm{Err}_1-\mathrm{Err}_2
=G_n-n^{-1}\|P_\Delta\varepsilon\|_2^2,
\]
and the final norm divided by \(\sigma^2\) is \(\chi_1^2\).

\subsection{Proof of Proposition~\ref{prop:rollout}}
Subtracting the true and fitted affine recursions gives
\[
e_{t+1}=(1-\widehat\tau_t)e_t+
(1-\bar y_t)(\widehat D_t-D_t)-
\bar y_t(\widehat R_t-R_t).
\]
Iteration from \(e_0=0\) yields Equation~\eqref{eq:rollout}. If
\(0\le\widehat\tau_t\le1\), every product has magnitude at most one. If also
\(\widehat\tau_t\ge\tau_{\min}>0\), the geometric series gives
\(|e_H|\le\max_k|\delta_k|\min\{H,1/\tau_{\min}\}\).

\section{SOURCE-TO-TARGET FINITE-SAMPLE BENCHMARK}

For completeness, suppose \((\alpha,\tau)\) is invariant across source
\(P\) and target \(Q\), and is estimated by homoskedastic least squares from
\(n_P\) independent source observations. In the source-centered basis, the
expected target excess one-step risk due only to estimation is
\begin{equation}
\frac{\sigma^2}{n_P}
\left[1+\frac{v_Q+(\mu_Q-\mu_P)^2}{v_P}\right].
\label{eq:transfercost}
\end{equation}
The first term is net-drift estimation; the second is turnover estimation plus
source-to-target leverage. A collapsed one-flow projection may additionally
change with the source state distribution. This benchmark shows why event
holdout can amplify either the one-flow projection cost or the two-flow
estimation cost; it gives no theorem that event holdout must favor two flows.

\section{EXPERIMENT COMPLETION RULE}

The main table is frozen before any model extension. The report must include all
eleven \(d_e^{\mathrm{step}}\) and \(d_e^{24}\) values, all stated event-level
uncertainty summaries, optimization diagnostics, HGB and persistence results,
and a clear H1/H2 decision. Only after that report is archived may a single
follow-up ablation be chosen, and its choice must follow the observed failure
mode rather than initiate a broad search.

\end{document}
